% Generated by Listas/make_statements.py from pratica-lecture-03.md, item 4. Do not edit.
\textbf{Source.} Tadelis (2013), \emph{Game Theory: An Introduction}, Exercise 7.2 ``Strategies and Equilibrium'', p.~147 {[}PDF 166{]}. Reproduced verbatim. Answer in the source: none in the book (the preface, p.~xiii, says about half are solved on Princeton UP's website; not checked).



\begin{quote}

\textbf{Strategies and Equilibrium:} Consider a two-player game in which player 1 can choose \(A\) or \(B\). The game ends if he chooses \(A\) while it continues to player 2 if he chooses \(B\). Player 2 can then choose \(C\) or \(D\), with the game ending after \(C\) and continuing again with player 1 after \(D\). Player 1 can then choose \(E\) or \(F\), and the game ends after each of these choices.



\begin{enumerate}

\def\labelenumi{\alph{enumi}.}

\item

  Model this as an extensive-form game tree. Is it a game of perfect or imperfect information?

\item

  How many terminal nodes does the game have? How many information sets?

\item

  How many pure strategies does each player have?

\item

  Imagine that the payoffs following choice \(A\) by player 1 are \((2, 0)\), those following \(C\) by player 2 are \((3, 1)\), those following \(E\) by player 1 are \((0, 0)\), and those following \(F\) by player 1 are \((1, 2)\). What are the Nash equilibria of this game? Does one strike you as more ``appealing'' than the other? If so, explain why.

\end{enumerate}

\end{quote}

